\section{Related Work}
\label{sec:related-work}

Two literatures bear on this paper, and they answer different questions. One asks whether the
\emph{setting} is real: is precision aquaculture a live engineering problem, and is edge
deployment of vision models on Apple Silicon or similar hardware a solved-enough practice to build
on. The other asks whether the \emph{safety pattern} is new: is putting a deterministic layer in
final authority over a probabilistic one a novel idea. The honest answer to the second question is
no, and stating that plainly, with the specific prior art named, is more useful to a reader than
pretending otherwise.

\subsection{Precision aquaculture and edge vision: an active, competitive setting}

Computer-vision monitoring of farmed fish is a mature enough sub-field to have its own taxonomy.
\citet{cui2025fishtracking}'s survey organizes the literature into detection, counting, behavior
classification, and activity analysis, and reports contemporaneous RAS-specific YOLO variants
reaching mAP50 above 0.90 -- higher than the 0.840 this paper's own detector achieves on a
510-image, single-class dataset. S.U.R.E.'s vision pipeline is not a vision-accuracy contribution
against that baseline, and this paper does not claim it is one. What the field's own review
literature does flag as missing is the closure from a detection to an actionable, safety-graded
decision: \citet{fitzgerald2025machinevision} state directly that most deployed vision systems in
aquaculture welfare monitoring never complete that trace. Precision-aquaculture reviews describe
detectors; they describe welfare taxonomies decomposed into input- and outcome-based indicators
\citep{frontiers2023welfare,pmc2021myfishcheck}; they do not, in general, describe a system that
follows one detector's confidence-threshold choice all the way to the single rule that consumes it.
Section~\ref{sec:results} traces exactly that path for one rule (\texttt{fish\_count==0}), and the
sturgeon-specific prior work is thin enough -- \citet{sturgeon2024cmscn}'s cross-modal stress
classifier is the closest, and it is a lab study, not a deployed system -- that a working,
bench-evaluated, sturgeon-specific decision pipeline is itself a modest but genuine contribution to that
narrow corner of the literature, independent of anything claimed about the LLM layer.

The motivating urgency claim -- that a dissolved-oxygen crash can kill stock within hours -- is not
this paper's own measurement; it is grounded in water-quality literature reporting species- and
temperature-dependent lethal thresholds in the 4.1--7.0~mg/L range
\citep{lindholm2023waterquality,agmrc2026waterquality}, and it should be read with that
species/temperature caveat attached rather than as a universal constant. On the edge-deployment
side, CoreML-on-Apple-Neural-Engine benchmarking is itself vendor-documented practice, not a novel
technique \citep{ultralytics2026coreml}, and the INT8 accuracy band this paper reports
(\(-0.0082\) mAP50) sits comfortably inside the 0.5--2.0-point range already reported for
calibrated quantization in the TensorRT/ONNX literature \citep{ultralytics2025tensorrt,
arxiv2025quantrobustness}. Community-documented practitioner reports of ONNX and TorchScript
export losing accuracy relative to eager PyTorch execution, independent of quantization, already
circulate informally on GitHub issue threads \citep{github2022yolov5issue,github2024onnxissue};
this paper's contribution there (Section~\ref{sec:results}) is quantifying that suspected mechanism
under one controlled, six-configuration sweep against a single held-out validation set, not
discovering it. None of this establishes novelty. It establishes that the ground the system stands
on -- fish detection, tracking with ByteTrack \citep{zhang2022bytetrack}, and export-format
benchmarking -- is real, current, and reasonably well understood, and that S.U.R.E. sits inside it
rather than reinventing it. Deployment context matters here too: on-device inference is favored
specifically because sub-100ms latency and unreliable rural connectivity make cloud round-trip
unattractive \citep{roboflow2025cloudedge}; this paper's latency numbers measure model inference
only, and should not be compared directly to reported end-to-end field latencies for a
cloud-served alternative.

\subsection{What this paper does not claim to have invented}

The pivot matters enough to state on its own. This paper does \emph{not} claim to have invented
deterministic override of a probabilistic component, interlocks, or runtime shielding. That
concept is decades old in industrial safety engineering: Safety Instrumented Systems, formalized
under IEC 61508 and IEC 61511, already specify a diagnosed, non-bypassable final-authority layer
for enumerable hazards, with the entire discipline organized around the assumption that the
protected process cannot be trusted to self-correct \citep{aiche2009beacon,icheme2026iec61511}.
The reinforcement-learning safety literature arrives at the same idea from a different direction:
a \emph{shield} vetoes or corrects a learned policy's proposed action before it reaches the
environment, and that literature independently states the same boundary condition this paper's
Formalization~\ref{form:enumerable} derives from first-hand observation of the system -- shielding
requires a formally specifiable, enumerable safety condition, and offers no guarantee once the
thing being checked is open-ended \citep{arxiv2024verifshielding,arxiv2025watchdogs}. Practitioner
writing on LLM guardrails has, by 2025--2026, already converged on almost the same slogan this
paper's system embodies -- the model proposes, a rule engine disposes
\citep{rulebricks2026deterministic,danilec2026deterministic,arthur2025guardrails,shields2024acm} --
so even the specific framing is not new currency.

One architectural distinction is worth naming rather than glossing over, because it is the one
place this paper's independence claim is weaker than the industrial-safety literature it borrows
language from. IEC 61511's Independent Protection Layer criteria are written around \emph{physical}
independence: separate sensors, separate actuators, separate failure paths. S.U.R.E.'s
independence is \emph{logical}, at the code level -- \texttt{backend/rules.py} is a single Python
module imported by both the system's decision path and the evaluation harness, and
Section~\ref{sec:system-architecture} treats that shared-import property as the thing doing the
safety work. A shared module that both paths import is a strictly weaker independence guarantee
than physically separate hardware, and this paper does not pretend otherwise; it is, however, a
real and verifiable property (Section~\ref{sec:results} reports a direct object-identity check,
not a code-review claim), and it is the property that actually failed once in this system's own
history -- an evaluation harness that had silently drifted from the system's real decision logic, on exactly the
scenario (\texttt{fish\_count==0}) this paper's vision-safety trace revisits.

\subsection{LLM guardrails and agentic reliability: a critical read, not a bibliography}

Most published guardrail architectures for LLMs re-validate a model's output with a second model,
a schema check, or a human reviewer -- not a dependency-free deterministic rule engine whose
escalation-only precedence is enforced by literally sharing code with the system under test. That
gap is worth naming before this paper's own contribution, because it is the specific thing most of
this cluster does not do. Academic proposals for LLM-mediated industrial control already exist at
the framework or simulation level: \citet{arxiv2024industrialllm} and a related line of work on
agentic industrial control \citep{arxiv2025agenticindustrial} propose exactly the pattern of
routing every LLM proposal through an external validator before actuation. What distinguishes
S.U.R.E. from that closest same-domain-class prior art is not the pattern but the evidentiary
status: those are proposed or simulated architectures, and this is an implemented system, evaluated
offline against synthetically generated sensor data, with a documented near-miss and eleven
independently verified experiments behind its numbers. Clinical
LLM safety offers a useful contrast in the opposite direction. There, the dominant backstop is
human-in-the-loop review, and residual hallucination rates remain non-trivial even under
mitigation -- one adversarial-robustness study reports failure rates near 44\% for clinical
decision support under adversarial prompting \citep{commsmed2025hallucination,bmc2026hallucination}.
A human reviewer does not scale to a 24-hour, unattended aquaculture deployment; substituting an
always-available, unit-tested deterministic module for that human is the design choice this paper
measures the consequences of, not a claim that the substitution itself is conceptually new.

On the LLM component's own reliability, this paper's negative result for agentic tool-calling
(Section~\ref{sec:results}) sits inside, rather than against, a growing and unresolved literature:
small on-device models' tool-calling reliability is described as an actively contested research
area with mixed findings across studies \citep{arxiv2510slmagentic}, and independent 2025--2026
benchmarks of comparably sized models report similarly severe failure rates
\citep{devto2026llama3b,arxiv2606toolenv}. That does not make AQUA-1B's 0\%/0\% result
uninteresting; it makes it one more literature-consistent data point rather than an outlier
demanding a special explanation. The constant-answer artifact this paper flags in AQUA-7B's
tool-selection benchmark is a named, established concern in benchmark-evaluation research more
broadly -- flagging near-zero-entropy answer distributions as a scoring artifact rather than
evidence of discrimination is not new \citep{arxiv2024artifacts} -- and this paper's contribution
there is a narrow application of that established discipline to a new task class (closed-menu
agentic tool selection), not the invention of artifact-aware evaluation.

The paper's own fabrication finding -- a model reaching the categorically correct status label
through a free-text reasoning chain that cites a sensor value never present in the input -- sits
closest to a strand of hallucination research that most benchmark work cannot reach, because most
benchmark work grades only the final answer. General hallucination and confabulation research in
natural-language generation is by now a mature survey area \citep{ji2023hallucinationsurvey}, and
a broader taxonomy has recently extended that framing specifically to agentic systems
\citep{arxiv2509agenthallucinationsurvey}. The closest direct match found for this paper's specific
observation -- fabrication surviving inside an output whose enumerable field is correct -- is a
2026 study on internal representations as indicators of hallucination in agentic tool selection
\citep{arxiv2601toolhallucination}, which shares the underlying concern but not this paper's
specific empirical instance: a small, self-documented-undertrained local model, in an implemented
safety-critical decision loop, where the fabrication is invisible to precisely the evaluation
regime -- status-only, answer-only grading -- that dominates the guardrail and agentic-benchmark
literature this cluster otherwise reviews. That is the gap this paper's Section~\ref{sec:results}
actually fills, and it is worth stating plainly rather than folding into a longer list: an
evaluation regime built to check an enumerable field will not catch a wrong reason attached to a
right decision, because it was never designed to look.

\subsection{Retrieval calibration and drift-gated MLOps: careful application, not invention}

The two remaining clusters ground this paper's RAG and MLOps components as applications of
active, already-established sub-literatures, and this paper treats them exactly that way.
Asymmetric retrieval training explains, mechanically, why the deployed e5-small embedding model
outperforms a symmetric sentence-similarity model as passage length grows -- the query/passage
prefix scheme \citet{wang2022e5} introduce is the reason, not a novel discovery of this paper.
Per-corpus similarity-threshold calibration is itself an established discipline, not a one-off
choice: recent work argues explicitly against a universal similarity threshold in favor of
corpus-specific precision/recall calibration \citep{researchgate2025ragprecision,
braintrust2025ragmetrics}, and this paper's 0.85-over-0.84 threshold selection
(Section~\ref{sec:results}) is exactly that discipline applied, with the asymmetric cost stated
rather than assumed. It is also bounded by a limit this paper does not try to talk around: recent
formal work on embedding-similarity-based hallucination gating shows there are certified limits to
how much a similarity threshold alone can prevent \citep{arxiv2026semanticillusion,
arxiv2025haltrag}, and the overlapping positive/negative similarity distributions this paper
measures on its own small corpus independently confirm that caution -- a threshold, however
carefully chosen, is a mitigation, not an elimination, of retrieval-driven fabrication risk.

Population Stability Index drift detection has the same status: the pitfall this paper's own
MLOps sweep rediscovers -- that binning choice materially changes PSI's sensitivity, and that
naive equal-width binning can behave badly under a narrow-banded score distribution -- is already
documented, standard practice in the credit-scoring literature PSI originates from
\citep{arize2026psi,arize2026binning,fiddler2026psi,yurdakul2018psi}. Champion-challenger
promotion gating with a minimum-improvement margin before shipping a retrained model is likewise
widely documented MLOps practice \citep{datarobot2026champion,snowflake2026champion}, as is the
broader tension between scheduled and drift-triggered retraining
\citep{wwt2026mlopsdrift,researchgate2025autodrift}. What is not already documented, as far as this
review can establish, is the specific combination this system runs: PSI over YOLO detection-
confidence deciles, narrow-banded for a structurally different reason than a credit score is
narrow-banded, feeding a champion-challenger gate for a vision-detection registry rather than a
tabular classifier. That is a real but modest contribution -- an application of two known
techniques to a new signal type and a new model class -- and this paper states it at that scale,
not as a discovery of either PSI or champion-challenger gating.

\subsection{What is actually new here}

After three clusters of prior art, the honest novelty locus is narrow and specific: the same
narrow-enumerable-interface, escalation-only discipline, independently measured across five
architecturally distinct layers of one implemented system -- decision fusion, agentic tool
routing, retrieval, MLOps retraining, and edge vision -- including one full disclosed negative
result and a fabrication-detection finding the Safety-Instrumented-System and shielding
literatures have no equivalent for, because those literatures are built around enumerable
action or state spaces, not free-text reasoning that can be factually wrong while the enumerable
output next to it is correct. Section~\ref{sec:results} is the actual argument for this claim; this
section's job was only to establish that nothing else in the surrounding literature already makes
it.
